\documentclass[11pt]{article}

\usepackage[margin=1in]{geometry}
\usepackage[T1]{fontenc}
\usepackage{lmodern}
\usepackage{microtype}
\usepackage{booktabs}
\usepackage{tabularx}
\usepackage{array}
\usepackage{enumitem}
\usepackage{xcolor}
\usepackage{xurl}
\usepackage{hyperref}

\hypersetup{
  colorlinks=true,
  linkcolor=black,
  urlcolor=blue,
  pdftitle={Technical Report: Provided-Text Conditioned Audio Delivery Ranking},
  pdfauthor={WhereIsAI-Lingnan}
}

\newcolumntype{Y}{>{\raggedright\arraybackslash}X}
\newcommand{\hashline}[2]{\shortstack[l]{\texttt{#1}\\\texttt{#2}}}
\setlist[itemize]{leftmargin=1.4em,itemsep=0.2em,topsep=0.25em}
\setlength{\parskip}{0.35em}
\setlength{\parindent}{0pt}
\sloppy

\title{\textbf{Technical Report: Provided-Text Conditioned\\Audio Delivery Ranking for EmpathyEval Track 1}}
\author{WhereIsAI-Lingnan}
\date{Final Release, July 2026}

\begin{document}
\maketitle

\begin{abstract}
This report documents the final Phase 1 submission system,
\texttt{provided\_\allowbreak text\_\allowbreak conditioned\_\allowbreak audio\_\allowbreak delivery\_\allowbreak ranking}. The system
uses a single \texttt{AudioDeliveryPairwiseMLP} trained on all official labeled
training data. Candidate decisions use prosody, Emotion2Vec, WavLM, task, and
text-length metadata only. ASR, teacher models, model fusion, fallback rules,
speech-quality features, and test-time training are excluded. The frozen
submission covers 530 questions and has been reproduced byte-for-byte from the
released final checkpoint and retained final feature inputs.
\end{abstract}

\section{Public Artifacts and Integrity}

The public code repository is
\url{https://github.com/Beichen-Hg/WhereIsAI-Lingnan}. The public Hugging Face
model repository is
\url{https://huggingface.co/BertramM/whereisai-lingnan-track1-audio}. The
Hugging Face repository contains only the final model release files:
\texttt{model.pt}, \texttt{config.yaml}, \texttt{feature\_group\_importance.json},
the model card, and a SHA-256 manifest. Competition data, cached feature files,
pretrained-model copies, and per-question test predictions are not publicly
released in those repositories.

\begin{table}[h]
\centering
\small
\begin{tabularx}{\linewidth}{@{}lY@{}}
\toprule
Artifact & SHA-256 \\
\midrule
Final checkpoint \texttt{model.pt} & \hashline{24303cb353b767c6f8c964876beec}{85cbf868633aa88e3f291cf8004a5e23083} \\
Checkpoint configuration \texttt{config.yaml} & \hashline{d26c614e9eb27207bd6c9acdc00419e}{64e998ec32899410e159564af1d9d5038} \\
Frozen submission JSONL & \hashline{f84f1a12568f2556db273fc3b70e7d87}{b6319e83cf086c7772405e3abd185222} \\
\bottomrule
\end{tabularx}
\caption{Final release integrity identifiers.}
\end{table}

\section{Final Method}

\subsection{Model}

The final learned model is \texttt{AudioDeliveryPairwiseMLP}. For each question,
the scorer constructs candidate-level representations and converts them into
pairwise features for ranking the available response audios. The checkpoint
contains the model state dictionary, the 1,775 feature names, and the fitted
feature normalization statistics needed by inference.

The exact feature groups are:

\begin{itemize}
  \item \textbf{Prosody}: basic acoustic delivery statistics extracted by the
  \texttt{librosa\_basic} backend.
  \item \textbf{Emotion2Vec}: audio emotion embeddings from
  \texttt{emotion2vec/emotion2vec\_plus\_large}.
  \item \textbf{WavLM}: speech embeddings from
  \texttt{microsoft/wavlm-base-plus}.
  \item \textbf{Task}: task-type indicators.
  \item \textbf{Text metadata}: lengths of official JSON-provided context,
  user, and response text.
\end{itemize}

Candidate semantic text is taken from the official JSON field
\texttt{json\_provided}. No ASR transcript is used to override provided text.
FunASR is only an optional loading backend for the Emotion2Vec emotion model;
it is not used as an automatic-speech-recognition decision module.

\subsection{Excluded Components}

The final inference path rejects or does not invoke ASR semantic replacement,
teacher calls, test labels, fusion models, fallback models, speech-quality
feature groups, or test-time optimization. Dependency and submission audits are
run as part of final inference.

\subsection{Design Rationale and Decision Flow}

The task is treated as a ranking problem rather than a response-generation
problem. Given a shared conversational situation, the system selects the audio
candidate whose delivery is most appropriate. Official text provides stable
conditioning and length metadata only; it is not rewritten, semantically
replaced by ASR, or scored by an LLM judge.

For each question, the system constructs user/candidate prosody, Emotion2Vec,
WavLM, task, and text-metadata features. It then creates every unordered pair
of candidates. The MLP estimates the probability that the left candidate is
preferred to the right candidate from left, right, signed-difference,
absolute-difference, and safe-ratio features. Each candidate accumulates its
pairwise win probabilities, the score is normalized by the number of opponents,
and the highest mean score is returned. This one mechanism handles both
two-candidate and three-candidate questions without a separate model or a
fusion rule.

Emotion2Vec and WavLM have complementary roles. Emotion2Vec contributes
emotion-oriented audio representations, while WavLM contributes general speech
representations. Prosody provides interpretable delivery statistics. The final
MLP learns their user-to-candidate relationships together with the official
task indicators and provided-text length metadata.

\section{Training Protocol and Parameters}

All official labeled training examples were used for the final refit. Phase 1
test examples were used only for blind manifest construction, feature
extraction, inference, and submission auditing. No Phase 1 test labels are
loaded by the final inference entry point.

\begin{table}[h]
\centering
\small
\begin{tabular}{@{}ll@{}}
\toprule
Setting & Final value \\
\midrule
Model & \texttt{AudioDeliveryPairwiseMLP} \\
Feature count & 1,775 \\
Hidden dimensions & [256, 128] \\
Dropout & 0.15 \\
Learning rate & 0.0003 \\
Weight decay & 0.001 \\
Batch size & 128 \\
Epochs & 6 \\
Random seed & 2026 \\
Training data & All official labeled training data \\
\bottomrule
\end{tabular}
\caption{Final refit configuration.}
\end{table}

The all-labeled training metrics are retained as training diagnostics and are
not claimed as held-out generalization results.

\section{Runtime Environment}

The final release was verified on Linux with Python 3.10.20, PyTorch
2.9.0+cu128, CUDA 12.8, NumPy 2.2.6, PyYAML 6.0.3, librosa 0.11.0, SoundFile
0.13.1, Transformers 5.8.1, Hugging Face Hub 1.14.0, FunASR 1.3.1, and
ModelScope 1.36.3. GPU execution accelerates feature extraction; the final
scorer also supports \texttt{DEVICE=cpu}. A compatible PyTorch wheel must be
selected for the target machine. Full audio-feature reconstruction may require
system audio tooling such as \texttt{ffmpeg}, depending on the selected
FunASR/torchaudio loading path.

\section{Reproduction Procedure}

The GitHub repository contains the final-only source code, scripts, tests, and
runtime dependency specification. Install dependencies and download the public
Hugging Face checkpoint to the expected project-relative model directory:

\begin{verbatim}
python -m pip install -r requirements.txt
cd artifacts/final_refit
hf download BertramM/whereisai-lingnan-track1-audio \
  --local-dir audio_delivery_pairwise_all_labeled_clean_no_quality_v1
cd ../..
\end{verbatim}

\subsection{Normal Inference}

With organizer-provided Phase 1 assets and the corresponding final feature files
available under \texttt{artifacts/}, generate predictions with:

\begin{verbatim}
PYTHON_BIN=python DEVICE=cpu bash scripts/142_infer_phase1_full_audio_delivery.sh
\end{verbatim}

\subsection{Full Feature Reconstruction}

When starting from organizer-provided raw Phase 1 assets, build the final
provided-text manifest and feature tables before normal inference:

\begin{verbatim}
PYTHON_BIN=python bash scripts/140_build_phase1_full_provided_text_manifest.sh
PYTHON_BIN=python bash scripts/141_build_phase1_full_provided_text_features.sh
PYTHON_BIN=python DEVICE=cpu bash scripts/142_infer_phase1_full_audio_delivery.sh
\end{verbatim}

\subsection{Exact Frozen-Submission Check}

When the local frozen audit file is available, the following command also
compares the output byte-for-byte with the frozen candidate:

\begin{verbatim}
PYTHON_BIN=python DEVICE=cpu bash scripts/155_reproduce_final_submission.sh
\end{verbatim}

The available scripts are:

\begin{itemize}
  \item \texttt{140\_build\_phase1\_full\_provided\_text\_manifest.sh}: build the
  provided-text Phase 1 manifest.
  \item \texttt{141\_build\_phase1\_full\_provided\_text\_features.sh}: extract
  final prosody, Emotion2Vec, and WavLM features and build the final feature
  table.
  \item \texttt{142\_infer\_phase1\_full\_audio\_delivery.sh}: execute final
  inference with the released checkpoint.
  \item \texttt{155\_reproduce\_final\_submission.sh}: run inference and compare
  the generated JSONL byte-for-byte with the frozen candidate when that local
  audit file is available.
\end{itemize}

% Keep the expected model path explicit for evaluators reading the PDF.
\begin{verbatim}
artifacts/final_refit/audio_delivery_pairwise_all_labeled_clean_no_quality_v1/
\end{verbatim}

\section{Validation and Submission Coverage}

The final code tests passed 15/15. Final inference passed both the submission
format audit and the dependency audit. The frozen submission contains 530 rows:

\begin{table}[h]
\centering
\small
\begin{tabular}{@{}lrr@{}}
\toprule
Source & Questions & Answer distribution \\
\midrule
GigaSpeech & 200 & A=119, B=81 \\
MELD & 210 & A=86, B=62, C=62 \\
EmoV-DB & 120 & A=64, B=56 \\
\midrule
Total & 530 & A=269, B=199, C=62 \\
\bottomrule
\end{tabular}
\caption{Frozen Phase 1 submission coverage.}
\end{table}

\section{Submission Checklist}

\begin{itemize}
  \item The Hugging Face model repository is public and ungated.
  \item The released checkpoint and configuration hashes are recorded above.
  \item The public GitHub repository contains code, runtime specifications,
  tests, and this report.
  \item The competition form should contain the public Hugging Face URL and
  the captain name required by the organizers.
\end{itemize}

\end{document}
